% test-004.tex -- comando \spmoney
%
% Compilar:  lualatex-dev test-004.tex
% Validar:   show-pdf-tags --xml test-004.pdf | rnv-wrapp
%            verapdf --flavour ua2 --format text test-004.pdf
%
% Cada caso es un parrafo numerado, para poder referirse a el al
% escuchar el PDF. Los comentarios "doc:" indican lo que documenta
% el manual de usuario; en el resto, la lectura se revisa escuchando.
% Las entradas invalidas (que detienen la compilacion) no van aqui.
% La tabla completa de monedas y alias va en test-005.
\DocumentMetadata
  {
    lang=es-CL,
    pdfversion=2.0,
    pdfstandard=ua-2,
    tagging=on,
    tagging-setup={math/setup={mathml-SE}},
  }
\documentclass{article}
\usepackage[spanish]{babel}
\usepackage{unicode-math}
\usepackage{spintent}
\usepackage{hyperref}
\hypersetup
  {
    pdfdisplaydoctitle = true,
    pdftitle           = {test-004: spmoney},
  }
\begin{document}

\section*{A. Formas del argumento}

% doc: «500 pesos»
1. Solo la cifra: $\spmoney{500}$.

% doc: «500 dólares»
2. Cifra y alias: $\spmoney{500 dolares}$.

% doc: «500 pesos chilenos»
3. Alias ISO en mayúsculas: $\spmoney{500 CLP}$.

% doc: «500 pesos chilenos»
4. Alias ISO en minúsculas: $\spmoney{500 clp}$.

% doc: «500 euros»
5. Alias de otra moneda: $\spmoney{500 euro}$.

% doc: «500 euros»
6. Símbolo del euro: $\spmoney{500 €}$.

% doc: «500 libras»
7. Símbolo de la libra: $\spmoney{500 £}$.

% doc: «500 libras»
8. Alias de la libra: $\spmoney{500 libra}$.

% doc: «500 dólares estadounidenses»
9. Dólar en formato ISO mayúsculas: $\spmoney{500 USD}$.

% doc: «500 dólares estadounidenses»
10. Dólar en formato ISO minúsculas: $\spmoney{500 usd}$.

% doc: «menos 500 pesos»
11. Cifra negativa: $\spmoney{-500}$.

% doc: «1234 coma 56 pesos»
12. Cifra decimal: $\spmoney{1234,56}$.

% doc: «1234 coma 56 euros»
13. Cifra decimal con alias: $\spmoney{1234,56 euro}$.

% doc: «0 coma 5 pesos»
14. Cifra menor que uno: $\spmoney{0,5}$.

\section*{B. Singular, plural y millones}

% doc: «1 peso»
15. Uno: $\spmoney{1}$.

% doc: «1 euro»
16. Uno con alias: $\spmoney{1 euro}$.

% doc: «2 pesos»
17. Dos: $\spmoney{2}$.

% doc: «0 pesos»
18. Cero: $\spmoney{0}$.

% doc: «menos 1 peso»
19. Uno negativo: $\spmoney{-1}$.

% doc: «1 coma 5 pesos»
20. Uno con decimales: $\spmoney{1,5}$.

% doc: «1000 pesos»
21. Mil: $\spmoney{1000}$.

% doc: «1000000 de pesos»
22. Un millón exacto: $\spmoney{1000000}$.

% doc: «2000000 de pesos»
23. Dos millones exactos: $\spmoney{2000000}$.

% doc: «1000000 de dólares»
24. Un millón con alias: $\spmoney{1000000 dolares}$.

% doc: «1000000 de pesos chilenos»
25. Un millón con alias ISO: $\spmoney{1000000 CLP}$.

% doc: «1500000 pesos»
26. Un millón y medio: $\spmoney{1500000}$.

\section*{C. Llave cifras-sep}

% doc: «100000 pesos»
27. Por defecto: $\spmoney{100000}$.

% doc: «100000 pesos»
28. Con false: $\spmoney[cifras-sep=false]{100000}$.

% doc: «100000 pesos»
29. Con false y espacios del usuario: $\spmoney[cifras-sep=false]{10\,00\,00}$.

% doc: «1234567 coma 89 pesos»
30. Decimales por defecto: $\spmoney{1234567,89}$.

% doc: «1234567 coma 89 pesos»
31. Decimales con false: $\spmoney[cifras-sep=false]{1234567,89}$.

\section*{D. Llave currency}

% doc: «500 pesos»
32. Peso, el valor por defecto: $\spmoney[currency=peso]{500}$.

% doc: «500 euros»
33. Euro por alias: $\spmoney[currency=euro]{500}$.

% doc: «500 euros»
34. Euro por símbolo: $\spmoney[currency=€]{500}$.

% doc: «500 libras»
35. Libra por alias: $\spmoney[currency=libra]{500}$.

% doc: «500 yenes»
36. Yen por alias: $\spmoney[currency=yen]{500}$.

% doc: «500 pesos chilenos»
37. Peso chileno en formato ISO: $\spmoney[currency=CLP]{500}$.

% doc: «500 dólares estadounidenses»
38. Dólar en minúsculas: $\spmoney[currency=usd]{500}$.

% doc: «500 dólares estadounidenses»
39. Dólar en mayúsculas: $\spmoney[currency=USD]{500}$.

% doc: «500 libras»
40. Prioridad del argumento: $\spmoney[currency=euro]{500 libra}$.

% doc: «500 dólares estadounidenses»
41. Prioridad del argumento con ISO: $\spmoney[currency=libra]{500 USD}$.

\section*{E. Llave sub-units}

% doc: «100 dólares con 50 centavos»
42. Con true: $\spmoney[sub-units=true]{100,5 dolares}$.

% doc: «100 coma 5 dólares»
43. Con false: $\spmoney[sub-units=false]{100,5 dolares}$.

% doc: «100 euros con 25 céntimos»
44. Euro con true: $\spmoney[sub-units=true]{100,25 euro}$.

% doc: «100 dólares con 05 centavos»
45. Subunidad de una sola cifra: $\spmoney[sub-units=true]{100,05 dolares}$.

% doc: «1 dólar con 01 centavo»
46. Subunidad en singular: $\spmoney[sub-units=true]{1,01 dolares}$.

% doc: «100 dólares»
47. Sin decimales: $\spmoney[sub-units=true]{100 dolares}$.

% doc: «500 pesos con 75 centavos»
48. Peso por defecto: $\spmoney[sub-units=true]{500,75}$.

% doc: «100 libras con 50 peniques»
49. Libra: $\spmoney[sub-units=true]{100,5 libra}$.

% doc: «100 dólares estadounidenses con 50 centavos»
50. Dólar en formato ISO: $\spmoney[sub-units=true]{100,5 USD}$.

% doc: «menos 100 dólares con 50 centavos»
51. Cifra negativa: $\spmoney[sub-units=true]{-100,5 dolares}$.

\section*{F. Llave money-math}

% doc: «500 pesos»
52. Con true, el valor por defecto: $\spmoney[money-math=true]{500}$.

% doc: «500 pesos»
53. Con false: $\spmoney[money-math=false]{500}$.

% doc: «500 dólares»
54. Con false y alias de dólar: $\spmoney[money-math=false]{500 dolares}$.

% doc: «500 euros»
55. Con false y otra moneda: $\spmoney[money-math=false]{500 euro}$.

% el won no está en la fuente matemática por defecto, sí en la de texto
% doc: «500 wons»
56. Con false y el won: $\spmoney[money-math=false]{500 won}$.

% doc: «500 pesos chilenos»
57. Con false y código ISO, que no cambia: $\spmoney[money-math=false]{500 CLP}$.

% doc: «100 euros con 50 céntimos euro»
58. Con false y subunidades: $\spmoney[money-math=false,sub-units=true]{100,5 euro}$.

\section*{G. Combinaciones de llaves}

% doc: «1234 euros con 50 céntimos»
59. Euro con subunidades: $\spmoney[currency=euro,sub-units=true]{1234,5}$.

% doc: «1234 euros con 50 céntimos»
60. Euro, subunidades y sin separador de cifras: $\spmoney[currency=euro,sub-units=true,cifras-sep=false]{1234,5}$.

% doc: «100 dólares con 50 centavos»
61. Subunidades y money-math false: $\spmoney[sub-units=true,money-math=false]{100,5 dolares}$.

% doc: «menos 1500000 pesos chilenos»
62. Peso chileno negativo con millones: $\spmoney[currency=CLP]{-1500000}$.

% doc: «1000000 de libras con 25 peniques»
63. Libra con millones y subunidades: $\spmoney[currency=libra,sub-units=true,cifras-sep=false]{1000000,25}$.

\section*{H. Dentro de fórmulas}

% doc: «500 pesos más 250 pesos es igual a 750 pesos»
64. Suma e igualdad: $\spmoney{500} + \spmoney{250} = \spmoney{750}$.

% doc: «500 euros es menor que 600 dólares»
65. Monedas distintas: $\spmoney{500 euro} < \spmoney{600 dolares}$.

% doc: «1 dólar con 50 centavos más 2 dólares»
66. Llaves distintas en una misma fórmula: $\spmoney[sub-units=true]{1,5 dolares} + \spmoney{2 dolares}$.

% doc: «100 dólares con 50 centavos»
67. En fórmula destacada:
\[ \spmoney[sub-units=true]{100,5 dolares} \]

% doc: «la fracción con numerador 500 pesos y denominador 2 pesos»
68. En una fracción: $\frac{\spmoney{500}}{\spmoney{2}}$.

% doc: «p es igual a 1500 pesos chilenos»
69. Con una variable: $P = \spmoney{1500 CLP}$.

\section*{I. Alias con tilde y símbolo escapado}

% el alias se acepta con tilde, con mayúscula inicial y en singular
% doc: «500 dólares»
70. Alias con tilde: $\spmoney{500 dólares}$.

% doc: «1 dólar»
71. Alias con tilde en singular: $\spmoney{1 dólar}$.

% el símbolo del peso escrito como \$ equivale a $
% doc: «500 pesos»
72. Símbolo escapado: $\spmoney{500 \$}$.

% doc: «500 pesos»
73. Símbolo escapado con money-math false: $\spmoney[money-math=false]{500 \$}$.

% la llave currency admite lo mismo que el argumento
% doc: «500 pesos»
74. Llave currency con símbolo escapado: $\spmoney[currency=\$]{500}$.

% doc: «500 dólares»
75. Llave currency con tilde: $\spmoney[currency=dólares]{500}$.

% doc: «100 dólares con 50 centavos»
76. Alias con tilde y subunidades: $\spmoney[sub-units=true]{100,5 dólares}$.

% doc: «100 pesos con 50 centavos»
77. Símbolo escapado y subunidades: $\spmoney[sub-units=true]{100,5 \$}$.

% doc: «100000 dólares»
78. Alias con tilde y espacios matemáticos: $\spmoney[cifras-sep=false]{10\,00\,00 dólares}$.

\end{document}
